\\maketitle
\\begin{abstract}
    本文针对2023年CUMCM A题航空运输网络优化问题，建立了基于图优化的航班时刻表设计与调度模型。首先将航空网络抽象为图结构，节点表示机场，边表示航线；其次构建多目标优化模型，以最大化航班频次、最小化延误时间和降低运营成本为目标；最后采用改进的NSGA-II算法求解Pareto前沿。
    
    \\textbf{针对问题一}，我们建立了航空网络的图论模型。将$n$个机场抽象为图$G=(V,E)$的顶点，航线抽象为带权有向边，权重包括飞行时间、航班频次和运营成本。
    
    \\textbf{针对问题二}，我们构建了航班时刻表优化模型。引入整数变量$x_{i,j,t}$表示从机场$i$到$j$在时刻$t$是否开行航班，建立混合整数规划模型：
    $$\\max Z = \\sum_{i,j,t} r_{i,j} x_{i,j,t} - \\sum_{i,j,t} c_{i,j} x_{i,j,t}$$
    
    \\textbf{针对问题三}，我们设计了改进的NSGA-II算法。引入精英保留策略和自适应交叉变异概率，提高算法收敛速度和多样性。实验结果表明，算法在100个机场规模的网络上可在30分钟内求得高质量解。
    
    本研究为航空公司优化航班时刻表、提高运营效率提供了科学的决策支持。
    
    \\keywords{航空网络\quad 图优化\quad 混合整数规划\quad NSGA-II\quad 航班调度}
\\end{abstract}
